\chapter{Introducere}

\section{Motivație}

În contextul actual, când conștientizarea asupra importanței unei alimentații sănătoase este în continuă creștere, capacitatea de a înțelege și interpreta corect informațiile nutriționale de pe etichetele alimentelor devine esențială. Această lucrare își propune să dezvolte o soluție tehnologică care să simplifice procesul de analiză a etichetelor nutriționale prin utilizarea tehnologiilor de recunoaștere optică a caracterelor (OCR) și a inteligenței artificiale.

Motivația pentru dezvoltarea acestei aplicații provine din experiența personală acumulată în domeniul culturismului și nutriției. În acest context, înțelegerea precisă a compoziției alimentelor nu este doar o preferință, ci o necesitate. Multe produse alimentare conțin liste complexe de ingrediente, cu denumiri științifice sau tehnice care sunt greu de înțeles pentru consumatorul obișnuit. Aceste ingrediente sunt adesea deliberat denumite folosind terminologia științifică pentru a masca natura lor sau pentru a face lista de ingrediente mai puțin accesibilă publicului larg.

Din experiența personală, a fost necesar în repetate rânduri să explic altor persoane de ce anumite ingrediente sunt benefice sau dăunătoare, deoarece acestea nu reușeau să înțeleagă compoziția produselor pe care le consumau. De asemenea, chiar și pentru o persoană cu experiență în domeniu, interpretarea unor ingrediente necunoscute necesită cercetare suplimentară, ceea ce face procesul de selecție a alimentelor ineficient și consumator de timp.

\section{Problema identificată}

Consumatorii se confruntă cu multiple dificultăți în înțelegerea etichetelor nutriționale, o problemă documentată extensiv în literatura de specialitate. Studiile arată că, deși majoritatea consumatorilor declară că citesc etichetele, mulți nu le înțeleg pe deplin, ceea ce le limitează capacitatea de a face alegeri alimentare informate \cite{Campos2011, Grunert2010}. Această discrepanță între intenție și înțelegere subliniază nevoia de instrumente care să faciliteze interpretarea datelor.

Principalele obstacole identificate sunt:

\textbf{Complexitatea terminologiei}: Ingredientele sunt frecvent listate folosind denumiri științifice în loc de nume comune (de exemplu, acid ascorbic în loc de vitamina C, sau tocoferoli în loc de vitamina E), creând o barieră semnificativă pentru consumatorii fără cunoștințe de specialitate \cite{Grunert2010}.

\textbf{Lungimea listelor de ingrediente}: Produsele moderne, în special cele ultra-procesate, conțin adesea liste foarte lungi de ingrediente, pe care consumatorii nu au timpul sau răbdarea să le analizeze în detaliu în timpul cumpărăturilor.

\textbf{Ingrediente nefamiliare}: Multe produse conțin ingrediente precum aditivi, emulgatori sau stabilizatori, ale căror nume și funcții sunt necunoscute publicului larg, ceea ce poate genera suspiciune sau confuzie \cite{Machin2019}.

\textbf{Lipsa contextului nutrițional}: Chiar dacă un consumator recunoaște un ingredient, acesta poate să nu înțeleagă impactul său real asupra sănătății (de exemplu, cantitatea în care devine relevant) sau să aibă informații incomplete sau eronate, adesea influențate de dezinformarea online.

\section{Importanța înțelegerii informațiilor nutriționale}

Înțelegerea corectă a ingredientelor și a informațiilor nutriționale este crucială din mai multe perspective, având un impact direct asupra sănătății individuale și colective.

\textbf{Sănătatea personală}: O mai bună înțelegere a etichetelor este asociată cu alegeri alimentare mai sănătoase, cum ar fi un consum redus de grăsimi, sodiu și zahăr, și un consum crescut de fibre \cite{Campos2011}. Acest lucru contribuie la prevenirea bolilor cronice precum obezitatea, diabetul de tip 2 și bolile cardiovasculare.

\textbf{Identificarea gradului de procesare}: Abilitatea de a diferenția între alimentele naturale, minim procesate și cele ultra-procesate este esențială. Consumul ridicat de alimente ultra-procesate este legat de riscuri crescute pentru sănătate, însă consumatorii întâmpină dificultăți în a le identifica corect doar pe baza listei de ingrediente \cite{Machin2019}.

\textbf{Gestionarea alergiilor și intoleranțelor}: Pentru persoanele cu alergii sau intoleranțe alimentare, citirea corectă și completă a listei de ingrediente este o problemă de siguranță, prevenind reacții adverse potențial severe.

\textbf{Demitizarea ingredientelor}: O analiză obiectivă, bazată pe dovezi, ajută la combaterea dezinformării. Multe ingrediente considerate „periculoase" de publicul larg sunt, în realitate, inofensive sau chiar benefice în cantitățile utilizate, în timp ce altele, considerate sigure, pot avea efecte negative în contextul unei diete dezechilibrate. O aplicație informativă poate oferi claritate și poate promova decizii bazate pe știință.

\section{Obiectivele aplicației}

Aplicația dezvoltată în cadrul acestei lucrări își propune să ofere utilizatorilor un portal simplu și eficient pentru gestionarea și analiza informațiilor nutriționale. Obiectivele principale includ:

\textbf{Digitizarea automată a informațiilor}: Utilizarea tehnologiei OCR pentru extragerea automată a informațiilor nutriționale și a listelor de ingrediente din fotografii, eliminând necesitatea introducerii manuale a datelor.

\textbf{Interpretarea inteligentă a ingredientelor}: Implementarea unui sistem bazat pe inteligență artificială care să traducă listele complexe de ingrediente într-un format ușor de înțeles, explicând rolul și impactul fiecărui component.

\textbf{Evaluarea gradului de procesare}: Calcularea unui scor de procesare pentru fiecare aliment, bazat pe sisteme de clasificare recunoscute precum NOVA, ajutând utilizatorii să identifice rapid produsele mai naturale față de cele ultra-procesate.

\textbf{Funcționalități avansate de filtrare și comparare}: Oferirea de instrumente pentru filtrarea alimentelor după diverse criterii nutriționale (conținut de proteine, sodiu, zahăruri etc.) și posibilitatea de a compara produse similare.

\textbf{Bază de date colaborativă}: Crearea unei platforme unde utilizatorii pot contribui cu informații despre alimente, beneficiind astfel de o bază de date în continuă expansiune.

\section{Limitări și considerații etice}

\subsection{Acuratețea informațiilor generate de AI}

Aplicația utilizează modele de limbaj mari pentru interpretarea ingredientelor, ceea ce prezintă următoarele limitări:

\begin{itemize}
    \item \textbf{Posibilitatea de erori}: Modelele AI pot genera informații incorecte despre ingrediente sau efectele acestora asupra sănătății.
    \item \textbf{Lipsa validării medicale}: Informațiile generate nu sunt verificate de profesioniști în domeniul sănătății și nu trebuie să înlocuiască sfatul medical.
    \item \textbf{Actualizarea cunoștințelor}: Modelele pot avea informații învechite despre cercetările nutriționale recente.
\end{itemize}

\subsection{Măsuri de siguranță implementate}

Pentru a atenua aceste riscuri, aplicația implementează:
\begin{itemize}
    \item Disclaimere clare în interfață, care subliniază că aplicația este un instrument informativ și nu unul medical.
    \item Încurajarea constantă a utilizatorilor să consulte specialiști (medici, dieteticieni) pentru decizii legate de sănătate.
\end{itemize}